\BeleskeID{04-s002}
% element: 04-s002:e01 — Hindu/Indo-Arabik i istorijski razvoj: Aryabhata, Brahmagupta, al-Uqlidisi
% element: 04-s002:e02 — Precrtana digresija: vavilonski, rimski, majanski sistemi
% element: 04-s002:e03 — Radiks r i celobrojna osnova; skup cifara ci i za sada pozitivna vrednost V
% element: 04-s002:e04 — Pozicioni zapis, n celih i k razlomljenih mesta i tačka osnove
% element: 04-s002:e05 — Težinska suma V i težina ri
% element: 04-s002:d01

Indeks $i=0$ pripada poslednjoj cifri ispred tačke: njena težina je $r^0=1$. Pomeranje za jedno mesto ulevo množi težinu sa $r$, a pomeranje udesno deli je sa $r$. Prva cifra iza tačke zato ima indeks $-1$ i težinu $1/r$; tačka razdvaja nenegativne od negativnih indeksa.

Osnova određuje i težine i dozvoljene cifre. U osnovi $7$ cifra $6$ je dozvoljena, ali cifra $7$ nije: sedam jedinica zahteva prelazak na sledeće mesto i zapis $10_7$. Isti zapis cifara može imati različitu vrednost u različitim osnovama.
